\documentclass[11pt]{article}
\usepackage[a4paper,margin=28mm]{geometry}
\usepackage{amsmath,amssymb,amsthm}
\usepackage[colorlinks=true,linkcolor=blue,urlcolor=blue,citecolor=blue]{hyperref}
\newtheorem{theorem}{Theorem}
\newtheorem{lemma}{Lemma}
\theoremstyle{remark}
\newtheorem*{remark}{Remark}
\newcommand{\Q}{\mathbb Q}
\newcommand{\Z}{\mathbb Z}
\newcommand{\F}{\mathbb F}
\newcommand{\R}{\mathbb R}
\DeclareMathOperator{\End}{End}
\DeclareMathOperator{\Gal}{Gal}
\DeclareMathOperator{\GL}{GL}
\DeclareMathOperator{\PGL}{PGL}
\DeclareMathOperator{\Aut}{Aut}
\DeclareMathOperator{\Lie}{Lie}
\title{No rational points of order \(22\) on geometrically split abelian surfaces}
\author{}
\date{21 October 2026}

\begin{document}
\maketitle

\begin{theorem}\label{thm:main}
Let \(A/\Q\) be a geometrically split abelian surface. Then \(A(\Q)\)
contains no point of order \(22\).

More precisely, if \(A(\Q)[11]\neq0\), then \(A\) is simple over
\(\Q\), and \(A_{\overline{\Q}}\) is isogenous to \(E_1\times E_2\)
for elliptic curves \(E_1,E_2\) that are not isogenous over
\(\overline{\Q}\).

In particular, neither \([22]\) nor \([2,22]\) occurs as the rational
torsion subgroup of a geometrically split genus-\(2\) Jacobian over
\(\Q\).
\end{theorem}

Here ``geometrically split'' means that \(A_{\overline{\Q}}\) is
isogenous to a product of elliptic curves. We impose no restriction on
the degree of this isogeny, on the integral endomorphism ring, or on
the field of definition of the endomorphisms. The theorem does not
require a principal polarization, and its proof uses no computer
calculation.

Throughout, \(G_K=\Gal(\overline{\Q}/K)\) for a number field
\(K\subset\overline{\Q}\). For a prime \(p\), we write
\(T=T_p(A)\) and \(V=V_p(A)=T\otimes_{\Z_p}\Q_p\), and
\(\chi_p:G_{\Q}\to\F_p^\times\) for the mod-\(p\) cyclotomic
character.

\section{Inputs}\label{sec:inputs}

We use the following results on elliptic curves.
\begin{enumerate}
\item[(a)] \emph{Mazur} \cite[Theorem~(8)]{Mazur}. The torsion
subgroup of \(E(\Q)\), for an elliptic curve \(E/\Q\), is
\(\Z/m\) with \(m\leq10\) or \(m=12\), or \(\Z/2\times\Z/2\nu\) with
\(\nu\leq4\). In particular, no elliptic curve over \(\Q\) has a
rational point of order \(11\).

\item[(b)] \emph{Kenku--Momose} \cite[Theorem~(0.1)]{KM}. Let \(N\) be
a product of powers of \(2,3,5,7,11,13\), and let \(m\mid N\). If
\(X_1(m,N)\) has positive genus and
\[
 (m,N)\notin\{(1,11),(1,13),(1,14),(1,15),(1,16),(1,18),(2,10),(2,12)\},
\]
then \(Y_1(m,N)(k)=\emptyset\) for every quadratic field \(k\). We apply
this with \((m,N)=(1,22)\). The curve \(X_1(22)\) has genus
\(1+180/12-20/2=6\): the group \(\pm\Gamma_1(22)\) has index \(180\)
in \(\mathrm{SL}_2(\Z)\), and \(X_1(22)\) has \(20\) cusps and no
elliptic points. Hence \emph{no elliptic curve over a quadratic field
has a rational point of order \(22\)}. This agrees with the complete
list of torsion groups over quadratic fields, as restated in
\cite[Theorem~2.1]{LFN}.

\item[(c)] \emph{Le Fourn--Najman} \cite[Theorem~1.1]{LFN}. A
\(\Q\)-curve is an elliptic curve over a number field that is
isogenous over \(\overline{\Q}\) to all of its
\(\Gal(\overline{\Q}/\Q)\)-conjugates; curves with complex
multiplication are allowed. If \(E\) is a \(\Q\)-curve over a
quadratic field \(K\), then \(E(K)\) has no point of order \(11\).
Indeed, no group in the list of \cite[Theorem~1.1]{LFN} has an element
of order \(11\). The key case is \cite[Proposition~3.2]{LFN}, which
excludes the torsion group \(\Z/11\); by \cite[Theorem~2.1]{LFN}, it is
the only torsion group over a quadratic field that contains a point of
order \(11\).

\item[(d)] \emph{Bourdon--Clark} \cite[Theorem~4.1]{BC}. Let
\(\mathcal O\) be an order of conductor \(f\) in an imaginary quadratic
field \(M\) (denoted \(K\) in \cite{BC}), and put \(w=\#\mathcal O^\times\). Let \(M(f)\) be the ring
class field of conductor \(f\). It equals \(M(j(E))\) for every elliptic
curve \(E\) with complex multiplication by \(\mathcal O\). By part (a)
of the theorem, there is a positive integer \(T(\mathcal O,1,N)\) with
the following property. For every \(d\geq1\), there exist a number field
\(F\supseteq M(f)\) with \([F:M(f)]=d\) and an
\(\mathcal O\)-CM elliptic curve \(E/F\) with an \(F\)-rational point
of order \(N\) if and only if \(T(\mathcal O,1,N)\mid d\). Parts (d)
and (e) compute this integer for a prime \(N=\ell\geq5\); we take
\(a=0\) and \(b=1\) there. Then
\(T(\mathcal O,1,\ell)=\widetilde T/w\), where
\(\widetilde T=\ell^2-1\) if \(\ell\nmid f\) and \(\ell\) is inert in
\(M\) (part (d)(i)). In each of the remaining cases (d)(ii)--(v),
\(\widetilde T=\ell-1\).

For \(\ell=11\), we have \(\widetilde T\in\{120,10\}\). If
\(w\neq2\), then \(\mathcal O=\Z[i]\) or \(\mathcal O=\Z[\zeta_3]\).
Then \(f=1\), and \(11\) is inert in \(M\), because
\(11\equiv3\pmod4\) and \(11\equiv2\pmod3\). These two orders give
\(T=120/4=30\) and \(T=120/6=20\). In all other cases \(w=2\) and
\(T\in\{60,5\}\). Every value is divisible by \(5\). Consequently:
\emph{if an elliptic curve \(E\) with complex multiplication by an
order in \(M\) is defined over a number field \(F\supseteq M(j(E))\)
and has an \(F\)-rational point of order \(11\), then
\(5\mid[F:M(j(E))]\).}
\end{enumerate}

We also use two standard facts about \(p\)-adic representations.
\begin{enumerate}
\item[(e)] \emph{Lattice independence.} Let \(T_1,T_2\) be stable
\(\Z_p\)-lattices in a finite-dimensional \(\Q_p\)-representation of
a group. Then \(T_1/pT_1\) and \(T_2/pT_2\) have the same composition
factors. Indeed, the characteristic polynomial of each group element
on \(T_i/pT_i\) is the reduction of its characteristic polynomial on
the rational representation, and the Brauer--Nesbitt theorem applies.
The statement also applies to isomorphic representations, after
transporting a lattice through the isomorphism.

\item[(f)] \emph{Determinant.} Every abelian variety over \(\Q\) has a
polarization defined over \(\Q\). It induces a nondegenerate
alternating \(G_{\Q}\)-equivariant pairing
\(V_p(A)\times V_p(A)\to\Q_p(1)\). Hence the determinant of
\(G_{\Q}\) acting on \(V_p(A)\) is the square of the \(p\)-adic
cyclotomic character. Therefore, for every \(G_{\Q}\)-stable lattice
\(T'\subset V_p(A)\), the determinant of \(G_{\Q}\) on \(T'/pT'\) is
\(\chi_p^2\). The kernel of \(\chi_p^2\) is
\(G_{\Q(\zeta_p+\zeta_p^{-1})}\). For \(p=11\), this field has degree
\(5\) over \(\Q\).
\end{enumerate}

\section{Coprime torsion in an elliptic isogeny class}\label{sec:coprime}

\begin{lemma}\label{lem:coprime}
Let \(E_1,E_2\) be \(K\)-isogenous elliptic curves over a number field
\(K\). Suppose that \(E_1(K)\) contains a point of prime order \(r\),
and \(E_2(K)\) contains a point of a distinct prime order \(s\). Then
their \(K\)-isogeny class contains an elliptic curve with a
\(K\)-rational point of order \(rs\).
\end{lemma}

\begin{proof}
Choose a \(K\)-isogeny \(f:E_1\to E_2\). Let \(N_s\) be the
\(s\)-primary subgroup scheme of \(\ker(f)\), and factor \(f\) as
\[
 E_1\xrightarrow{\ q\ }E'=E_1/N_s\xrightarrow{\ f'\ }E_2.
\]
The degree of \(q\) is a power of \(s\), so \(q\) maps the given point
of order \(r\) to a point of order \(r\) in \(E'(K)\). The degree of
\(f'\) is prime to \(s\), so \(f'\) induces a Galois-equivariant
isomorphism \(E'[s]\simeq E_2[s]\). Thus \(E'(K)\) also contains a
point of order \(s\). The sum of the two points has order \(rs\).
\end{proof}

\section{Elliptic subvarieties and the endomorphism field}\label{sec:subvarieties}

Let \(A/\Q\) be an abelian surface. Put
\(D=\End^0(A_{\overline{\Q}})=\End(A_{\overline{\Q}})\otimes\Q\), and
let \(L\) be the smallest field over which all endomorphisms of
\(A_{\overline{\Q}}\) are defined. It is a finite Galois extension of
\(\Q\). The group \(\Gamma=\Gal(L/\Q)\) acts faithfully on \(D\) by
\(\Q\)-algebra automorphisms, \(\varphi\mapsto\varphi^\sigma\). For
\(\sigma\in G_{\Q}\), \(\varphi\in\End(A_{\overline{\Q}})\), and
\(x\in A(\overline{\Q})\), we have
\(\sigma(\varphi(x))=\varphi^\sigma(\sigma(x))\).

\begin{lemma}\label{lem:subvarieties}
\begin{enumerate}
\item[(a)] Every abelian subvariety \(B\) of \(A_{\overline{\Q}}\) is
the image of an endomorphism \(\varepsilon\) of \(A_{\overline{\Q}}\)
with \(\varepsilon|_B=[n]\) for some \(n\geq1\). In particular,
\(B\) is defined over \(L\).
\item[(b)] If \(B_1,B_2\) are abelian subvarieties of
\(A_{\overline{\Q}}\), then every homomorphism \(B_1\to B_2\) over
\(\overline{\Q}\) is defined over \(L\).
\end{enumerate}
\end{lemma}

\begin{proof}
(a) By Poincar\'e's complete reducibility theorem, there is an abelian
subvariety \(B'\) with \(B+B'=A_{\overline{\Q}}\) and \(B\cap B'\)
finite. The addition map \(\mu:B\times B'\to A_{\overline{\Q}}\) is an
isogeny. Choose an isogeny \(\nu\) in the other direction and
\(n\geq1\) with \(\mu\circ\nu=[n]\) and \(\nu\circ\mu=[n]\). Put
\(\varepsilon=\iota_B\circ\mathrm{pr}_1\circ\nu\), where \(\iota_B\) is
the inclusion. For \(b\in B\) we get
\(\varepsilon(b)=\mathrm{pr}_1(\nu(\mu(b,0)))=nb\). Thus
\(\varepsilon(A_{\overline{\Q}})=B\). Since \(\varepsilon\) is
defined over \(L\), every \(\sigma\in G_L\) satisfies
\(\sigma(B)=\varepsilon^\sigma(A_{\overline{\Q}})=B\). By Galois
descent, \(B\) is defined over \(L\).

(b) Let \(\varepsilon_1\) be as in (a) for \(B_1\), and let
\(\iota_2\) be the inclusion of \(B_2\). For a homomorphism
\(\varphi:B_1\to B_2\), the endomorphism
\(\psi=\iota_2\circ\varphi\circ\varepsilon_1\) of
\(A_{\overline{\Q}}\) is fixed by \(G_L\). Since \(\iota_2\) and
\(\varepsilon_1\) are defined over \(L\), every \(\sigma\in G_L\)
satisfies \(\iota_2\circ\varphi^\sigma\circ\varepsilon_1=\psi^\sigma
=\iota_2\circ\varphi\circ\varepsilon_1\). As \(\varepsilon_1\) maps
onto \(B_1\) and \(\iota_2\) is injective, \(\varphi^\sigma=\varphi\).
\end{proof}

\paragraph{The case of an elliptic curve over \(\Q\).}
Suppose that \(P\in A(\Q)\) has order \(11\), and that \(A\) contains
an elliptic curve \(E\) defined over \(\Q\). Put \(F=A/E\), an
elliptic curve over \(\Q\). Either \(P\in E(\Q)\), or \(P\) maps to a
point of order \(11\) in \(F(\Q)\). Both possibilities contradict
Mazur's theorem. Hence \(A(\Q)[11]\neq0\) implies that \(A\) is
simple over \(\Q\).

\paragraph{Splitting types.}
Now let \(A\) be geometrically split, so that
\(A_{\overline{\Q}}\) is isogenous to \(E_1\times E_2\). Exactly one
of the following holds.
\begin{enumerate}
\item[(N)] \(E_1\) and \(E_2\) are not isogenous over \(\overline{\Q}\).
\item[(S)] \(E_1\) and \(E_2\) are isogenous to an elliptic curve \(E\)
without complex multiplication. Then \(D\simeq M_2(\Q)\).
\item[(CM)] \(E_1\) and \(E_2\) are isogenous to an elliptic curve
\(E\) with complex multiplication by an order in an imaginary
quadratic field \(M\). Then \(D\simeq M_2(M)\), with centre \(M\).
\end{enumerate}
In cases (S) and (CM), every elliptic subvariety and every elliptic
quotient of \(A_{\overline{\Q}}\) is isogenous to \(E\): a nonzero
homomorphism between such a curve and \(E^2\) gives a nonzero
homomorphism between that curve and \(E\).
Sections~\ref{sec:endomorphism-field}--\ref{sec:small} show that
\(A(\Q)[11]=0\) in case (S), and Section~\ref{sec:cm} shows the same in
case (CM). Section~\ref{sec:completion} treats case (N).

\section{The endomorphism field of a non-CM square}\label{sec:endomorphism-field}

Assume case (S) and fix an isomorphism \(D\simeq M_2(\Q)\). By the
Skolem--Noether theorem, each \(\gamma\in\Gamma\) acts on \(D\) as
conjugation by a matrix \(S_\gamma\in\GL_2(\Q)\), unique up to
\(\Q^\times\). Since the action is faithful, \(\gamma\mapsto[S_\gamma]\)
embeds \(\Gamma\) into \(\PGL_2(\Q)\).

\begin{lemma}\label{lem:finite-subgroups}
The group \(\Gamma\) is cyclic of order \(1,2,3,4\), or \(6\), or
dihedral of order \(4,6,8\), or \(12\). We write these groups as
\[
 C_1,\ C_2,\ C_3,\ C_4,\ C_6,\qquad D_4,\ D_6,\ D_8,\ D_{12},
\]
where the subscript denotes the group order. The group \(D_4\) is the
Klein four group.
\end{lemma}

\begin{proof}
A finite subgroup of \(\PGL_2(\R)\) fixes a point of the hyperbolic
plane, and its action on the tangent plane at that point embeds it in
\(O(2)\). Hence it is cyclic or dihedral.

Let \([S]\in\PGL_2(\Q)\) have finite order \(n>2\). If \(S\) had a
repeated eigenvalue \(\lambda\), then \(S=\lambda(I+N)\) with \(N\)
nilpotent and nonzero, and \(S^n=\lambda^n(I+nN)\) would not be scalar.
Hence \(S\) has distinct eigenvalues \(\lambda_1,\lambda_2\). The order
of \([S]\) equals the order of \(\zeta=\lambda_1/\lambda_2\), so
\(\zeta\) is a primitive \(n\)-th root of unity, and
\[
 \frac{(\operatorname{tr}S)^2}{\det S}=2+\zeta+\zeta^{-1}\in\Q.
\]
The number \(\zeta+\zeta^{-1}\) has degree \(\varphi(n)/2\) over \(\Q\),
so \(\varphi(n)=2\) and \(n\in\{3,4,6\}\). A cyclic group of order
\(n\), and a dihedral group of order \(2n\), contain elements of order
\(n\). This gives the list.
\end{proof}

If \(\Gamma\in\{C_3,C_4,C_6,D_6,D_8,D_{12}\}\), let
\(\langle\rho\rangle\) be its rotation subgroup, of order
\(n\in\{3,4,6\}\). For a cyclic group it is the whole group; for a
dihedral group it is the cyclic subgroup of index \(2\). Every element
of \(\Gamma\) conjugates \(\rho\) to \(\rho\) or to \(\rho^{-1}\).

From now on until the end of Section~\ref{sec:small}, set \(p=11\),
\(k=\Q_{11}\), \(\mathcal O=\Z_{11}\), \(T=T_{11}(A)\), and
\(V=V_{11}(A)\). The isomorphism \(D\simeq M_2(\Q)\) extends to
\(D\otimes_{\Q}k\simeq M_2(k)\), which acts on \(V\). For
\(\sigma\in G_{\Q}\) with image \(\gamma\in\Gamma\),
\(x\in D\otimes_{\Q}k\), and \(v\in V\), we have
\begin{equation}\label{eq:semilinear}
 \sigma(xv)=(S_\gamma xS_\gamma^{-1})\,\sigma(v).
\end{equation}

\section{A residual tensor lemma}\label{sec:tensor}

\begin{lemma}\label{lem:tensor}
Assume case (S), and suppose that \(A(\Q)\) contains a point \(P\) of
order \(11\). Suppose that there is an \(\mathcal O\)-lattice
\(\Lambda\subset k^2\) with the following two properties.
\begin{enumerate}
\item[(i)] For every \(\gamma\in\Gamma\), we have
\(S_\gamma\Lambda=\lambda\Lambda\) for some \(\lambda\in k^\times\).
\item[(ii)] Some \(\gamma_0\in\Gamma\) has a representative
\(S\in k^\times S_{\gamma_0}\) with \(S\Lambda=\Lambda\) whose
reduction \(\overline S\in\GL(\Lambda/11\Lambda)\) has no eigenvector.
\end{enumerate}
Then \(\Q(\zeta_{11}+\zeta_{11}^{-1})\subseteq L\). In particular,
\(5\) divides \([L:\Q]\).
\end{lemma}

\begin{proof}
Let \(\mathcal R=\End_{\mathcal O}(\Lambda)\simeq M_2(\mathcal O)\),
a maximal order in \(M_2(k)\). By (i),
\(S_\gamma\mathcal R S_\gamma^{-1}=\End_{\mathcal O}(S_\gamma\Lambda)
=\mathcal R\) for every \(\gamma\). Put \(T'=\mathcal RT\). This is a
finitely generated \(\mathcal O\)-submodule of \(V\) containing \(T\),
hence a full lattice, and it is \(G_{\Q}\)-stable by
\eqref{eq:semilinear}.

Let \(e_{ij}\) be the matrix units of \(\mathcal R\simeq M_2(\mathcal O)\)
with respect to a basis of \(\Lambda\), and put \(W_0=e_{11}T'\). Then
\(T'=e_{11}T'\oplus e_{22}T'\), and \(e_{21}:e_{11}T'\to e_{22}T'\) is
an isomorphism with inverse \(e_{12}\). Hence
\(T'\simeq\Lambda\otimes_{\mathcal O}W_0\) as \(\mathcal R\)-modules,
with \(W_0\) free of rank \(2\). Reducing modulo \(11\), we obtain
\[
 \overline V=T'/11T'\simeq U\otimes_{\F_{11}}\overline W_0,
 \qquad U=\Lambda/11\Lambda,
\]
where \(\mathcal R/11\mathcal R=\End(U)\) acts on the first factor.

By (i), each \(\sigma\in G_{\Q}\), with image \(\gamma\in\Gamma\), has
a representative \(S_\sigma\in k^\times S_\gamma\) with
\(S_\sigma\Lambda=\Lambda\). It is unique up to \(\mathcal O^\times\),
so its reduction \(\overline S_\sigma\in\GL(U)\) is well defined up to
\(\F_{11}^\times\). By \eqref{eq:semilinear},
\(\sigma(xv)=(\overline S_\sigma x\overline S_\sigma^{-1})\sigma(v)\)
for \(x\in\End(U)\) and \(v\in\overline V\).

The modules \(\overline V\) and \(A[11]=T/11T\) are reductions of two
\(G_{\Q}\)-stable lattices in \(V\). By lattice independence, they have
the same composition factors. The point \(P\) spans a trivial
submodule of \(A[11]\), so \(\overline V\) has a trivial
\(G_{\Q}\)-composition factor.

Consider the group
\[
 \widetilde G=\{(\sigma,\beta)\in G_{\Q}\times\GL(U):
 \beta\in\F_{11}^\times\overline S_\sigma\},
\]
which acts on \(\overline V\) through \(\sigma\), and on \(U\) through
\(\beta\). Its composition factors on \(\overline V\) are the same as
those of \(G_{\Q}\). For \((\sigma,\beta)\in\widetilde G\), the map
\((\beta^{-1}\otimes1)\circ\sigma\) on \(\overline V\) commutes with
\(\End(U)\otimes1\). Hence it equals \(1\otimes w(\sigma,\beta)\) for a
unique \(w(\sigma,\beta)\in\GL(\overline W_0)\). Then \(\sigma\) acts
on \(\overline V\) as \(\beta\otimes w(\sigma,\beta)\), and \(w\) is a
homomorphism. Writing \(W\) for \(\overline W_0\) with this action, we
get an isomorphism of \(\widetilde G\)-modules
\[
 \overline V\simeq U\otimes W.
\]

The module \(U\) is irreducible. Indeed, for any \(\sigma_0\in G_{\Q}\)
with image \(\gamma_0\), the element \((\sigma_0,\overline S)\) lies in
\(\widetilde G\), and \(\overline S\) has no eigenvector by (ii). This
is the point where the projective Galois action on \(U\) is used, and
not only the structure of an order.

The module \(W\) is also irreducible. Otherwise its composition factors
are two characters \(\psi_1,\psi_2\), and the composition factors of
\(U\otimes W\) are the irreducible two-dimensional modules
\(U\otimes\psi_1\) and \(U\otimes\psi_2\). This contradicts the trivial
composition factor of \(\overline V\).

The subgroup \(H=\{(h,1):h\in G_L\}\) is normal in \(\widetilde G\) and
acts trivially on \(U\). The restriction \(W|_H\) is semisimple by
Clifford's theorem. Explicitly, the \(\widetilde G\)-translates of a
simple \(H\)-submodule of \(W\) are simple \(H\)-submodules, and their
sum is a nonzero \(\widetilde G\)-submodule, hence equal to \(W\). Now
\(\overline V|_H\simeq W|_H\oplus W|_H\) has a trivial composition
factor. Thus the semisimple module \(W|_H\) contains a trivial
submodule, so \(W^H\neq0\). Since \(H\) is normal, \(W^H\) is a
\(\widetilde G\)-submodule, so \(W^H=W\). Hence \(G_L\) acts trivially
on \(U\otimes W=\overline V\).

By fact (f) of Section~\ref{sec:inputs}, the determinant of
\(G_{\Q}\) on \(\overline V\) is \(\chi_{11}^2\). Therefore
\(\chi_{11}^2|_{G_L}=1\), and
\(\Q(\zeta_{11}+\zeta_{11}^{-1})\subseteq L\). This field has degree
\(5\).
\end{proof}

\section{Rotations of order 3, 4, and 6}\label{sec:rotations}

Assume case (S), let \(P\in A(\Q)\) have order \(11\), and suppose
that \(\Gamma\in\{C_3,C_4,C_6,D_6,D_8,D_{12}\}\), with rotation
\(\rho\) of order \(n\in\{3,4,6\}\). Choose a representative
\(R\in\GL_2(\Q)\) of \(\rho\). If \(\operatorname{tr}R=0\), the
Cayley--Hamilton theorem gives \(R^2=-\det R\), so \(\rho^2=1\).
Hence \(\operatorname{tr}R\neq0\), and we normalize \(R\) to have trace
\(1\). By the proof of Lemma~\ref{lem:finite-subgroups},
\(1/\det R=2+\zeta_n+\zeta_n^{-1}\), which equals \(1,2,3\) for
\(n=3,4,6\). Thus the characteristic polynomial of \(R\) is
\[
 X^2-X+\frac1c,\qquad c=1,2,3\quad\text{for }n=3,4,6,
\]
with discriminant \(-3\), \(-1\), \(-1/3\), respectively. Modulo
\(11\), these are \(8\), \(10\), and \(7\), since \(3^{-1}\equiv4\).
The nonzero squares modulo \(11\) are \(1,3,4,5,9\). Hence none of the
three discriminants is a square modulo \(11\). Since \(1/c\) is a unit,
the polynomial is irreducible modulo \(11\) with unit discriminant. So
\(k'=k[R]\) is the unramified quadratic extension of \(k\), and
\(\mathcal O[R]\) is its ring of integers.

\emph{Stable lattices.} The plane \(k^2\) is a one-dimensional
\(k'\)-vector space. An \(\mathcal O[R]\)-stable lattice is a finitely
generated \(\mathcal O[R]\)-submodule spanning it, hence of the form
\(\mathcal O[R]u\). Since \(11\) is a uniformizer of \(\mathcal O[R]\),
any two such lattices differ by a power of \(11\). Fix one, \(\Lambda\).

\emph{Property (i).} Let \(\gamma\in\Gamma\). Then
\(S_\gamma RS_\gamma^{-1}\) is a scalar multiple of \(R\) or of
\(R^{-1}=(1-R)/\det R\). It has trace \(1\), so it equals \(R\) or
\(1-R\). Hence \(S_\gamma\) normalizes \(\mathcal O[R]\). Thus
\(S_\gamma\Lambda\) is \(\mathcal O[R]\)-stable, and so it is a power of
\(11\) times \(\Lambda\).

\emph{Property (ii).} Take \(\gamma_0=\rho\) and \(S=R\). Then
\(R\Lambda\subseteq\Lambda\), and \(\det R=1/c\) is a unit, so
\(R\Lambda=\Lambda\). The reduction \(\overline R\) has the irreducible
characteristic polynomial \(X^2-X+1/c\) over \(\F_{11}\), so it has no
eigenvector.

Lemma~\ref{lem:tensor} now gives \(5\mid|\Gamma|\). This is impossible
for groups of order \(3,4,6,8,12\). Hence
\(\Gamma\in\{C_1,C_2,D_4\}\).

\section{The Klein four group}\label{sec:klein}

Assume case (S), let \(P\in A(\Q)\) have order \(11\), and suppose
that \(\Gamma=D_4\simeq C_2\times C_2\). The map
\[
 \delta:\Gamma\to\Z/2\Z,\qquad
 \gamma\mapsto v_{11}(\det S_\gamma)\bmod2,
\]
is well defined, because rescaling \(S_\gamma\) by \(t\) multiplies the
determinant by \(t^2\). It is a homomorphism. Since \(|\Gamma|=4\), its
kernel contains a nonidentity element \(r\). Choose a representative
\(R\in\GL_2(\Q)\) of \(r\). Since \(r^2=1\), we have \(R^2=aI\) with
\(a\in\Q^\times\). As \(R\) is not scalar, its minimal and
characteristic polynomials equal \(X^2-a\). Hence
\(\operatorname{tr}R=0\) and \(\det R=-a\), so \(v_{11}(a)\) is even.
Replacing \(R\) by \(11^{-v_{11}(a)/2}R\), we may assume that \(a\) is
an \(11\)-adic unit. Put
\[
 \mathcal B=\mathcal O[R]=\mathcal O+\mathcal OR.
\]

Since \(\Gamma\) is abelian, \(S_\gamma RS_\gamma^{-1}=\mu R\) with
\(\mu\in\Q^\times\) for every \(\gamma\in\Gamma\). Squaring gives
\(aI=\mu^2aI\), so \(\mu=\pm1\). We write \(\mu=\epsilon(\gamma)\).
This defines a homomorphism \(\epsilon:\Gamma\to\{\pm1\}\) with
\(\epsilon(r)=1\). In particular, every \(S_\gamma\) normalizes
\(\mathcal B\).

\paragraph{The nonsplit case.}
Suppose that \(a\) is not a square modulo \(11\). Then \(X^2-a\) is
irreducible modulo \(11\) with unit discriminant \(4a\), and
\(\mathcal B\) is the ring of integers of the unramified quadratic
extension of \(k\). The argument of Section~\ref{sec:rotations}
applies verbatim. All \(\mathcal B\)-stable lattices are powers of
\(11\) times one of them, \(\Lambda\). Each \(S_\gamma\) maps
\(\Lambda\) to a \(\mathcal B\)-stable lattice. The representative
\(R\) of \(r\) satisfies \(R\Lambda=\Lambda\), since \(\det R=-a\) is
a unit. Its reduction has the irreducible characteristic polynomial
\(X^2-\overline a\). Lemma~\ref{lem:tensor} gives \(5\mid4\), a
contradiction.

\paragraph{The split case.}
Suppose that \(a\) is a nonzero square modulo \(11\). By Hensel's
lemma, \(a=\alpha^2\) with \(\alpha\in\mathcal O^\times\). Put
\[
 e_1=\frac{\alpha+R}{2\alpha},\qquad e_2=\frac{\alpha-R}{2\alpha}.
\]
These are orthogonal idempotents in \(\mathcal B\) with
\(e_1+e_2=1\) and \(R=\alpha(e_1-e_2)\), so
\(\mathcal B\simeq\mathcal O\times\mathcal O\). In a basis of
\(k^2\) consisting of eigenvectors of \(R\), we have
\(e_1=E_{11}\) and \(e_2=E_{22}\). Put \(X=E_{12}\) and
\(Y=E_{21}\) in this basis. Then
\[
 XY=e_1,\qquad YX=e_2,\qquad e_1M_2(k)e_2=kX.
\]

\emph{The character \(\epsilon\) is nontrivial.} Otherwise every
\(S_\gamma\) commutes with \(R\), and is therefore diagonal in the
eigenbasis. The image of the diagonal torus in \(\PGL_2(k)\) is
isomorphic to \(k^\times\), via
\(\operatorname{diag}(x,y)\mapsto x/y\). The map
\(\Gamma\to\PGL_2(\Q)\to\PGL_2(k)\) is injective. Thus \(\Gamma\) would
embed in \(k^\times\). Finite subgroups of the multiplicative group of
a field are cyclic, but \(\Gamma\) is not.

Hence \(\ker\epsilon=\{1,r\}\). Let \(K=L^{\ker\epsilon}\), a quadratic
field. For \(\gamma\in\ker\epsilon\), \(S_\gamma\) commutes with
\(e_1\) and \(e_2\). For \(\gamma\notin\ker\epsilon\),
\(S_\gamma RS_\gamma^{-1}=-R\), so \(S_\gamma\) exchanges \(e_1\) and
\(e_2\).

\emph{Two blocks.} Put \(T'=\mathcal BT\). It is a full lattice in
\(V\), and it is \(G_{\Q}\)-stable by \eqref{eq:semilinear}. We have
\[
 T'=T_1\oplus T_2,\qquad T_i=e_iT'.
\]
By \eqref{eq:semilinear}, \(G_K\) preserves \(T_1\) and \(T_2\), and
every element of \(G_{\Q}\setminus G_K\) maps \(T_1\) onto \(T_2\) and
\(T_2\) onto \(T_1\). Put \(V_i=T_i\otimes_{\mathcal O}k=e_iV\). The
maps \(X:V_2\to V_1\) and \(Y:V_1\to V_2\) are inverse isomorphisms,
since \(YX=e_2\) and \(XY=e_1\). Thus each \(V_i\) has dimension
\(2\), and each \(T_i\) is free of rank \(2\).

\emph{The blocks differ by a character.} Let \(\sigma\in G_K\) have
image \(\gamma\in\ker\epsilon\). Since \(S_\gamma\) commutes with
\(e_1,e_2\), we have \(S_\gamma XS_\gamma^{-1}\in e_1M_2(k)e_2=kX\).
Write \(S_\gamma XS_\gamma^{-1}=\eta(\sigma)X\). This defines a
character \(\eta:G_K\to k^\times\), which factors through
\(\Gal(L/K)=\{1,r\}\). For \(r\) we compute
\[
 RXR^{-1}=(e_1-e_2)X(e_1-e_2)=-X,
\]
using \(X=e_1Xe_2\). Hence \(\eta\) is the quadratic character of
\(L/K\), with values \(\pm1\). It is trivial on \(G_L\), and its
reduction \(\overline\eta\) modulo \(11\) is nontrivial.

By \eqref{eq:semilinear}, \(\sigma(Xv)=\eta(\sigma)X\sigma(v)\) for
\(\sigma\in G_K\) and \(v\in V_2\). Thus \(X\) is a
\(G_K\)-isomorphism \(\eta\otimes V_2\to V_1\). Put
\(M_i=T_i/11T_i\). The lattice \(XT_2\subset V_1\) is
\(G_K\)-stable, and its reduction is isomorphic to
\(\overline\eta\otimes M_2\). By lattice independence in \(V_1\),
\[
 M_1^{\mathrm{ss}}\simeq\overline\eta\otimes M_2^{\mathrm{ss}}.
\]
This does not require \(X\) to carry \(T_2\) onto \(T_1\).

\emph{Both blocks contain the trivial character.} By lattice
independence, \(\overline V=T'/11T'\) has a trivial
\(G_{\Q}\)-composition factor, because \(A[11]\) does. As a
\(G_K\)-module, \(\overline V=M_1\oplus M_2\), so some \(M_i\) has a
trivial \(G_K\)-composition factor. Let
\(\tau\in G_{\Q}\setminus G_K\). Then \(t\mapsto\tau(t)\) maps
\(T_1\) onto \(T_2\), and \(\tau(\sigma t)=(\tau\sigma\tau^{-1})\tau(t)\).
So \(M_2\) is the \(\tau\)-conjugate of \(M_1\). A composition factor
\(\psi\) of one block gives the composition factor
\(\sigma\mapsto\psi(\tau^{-1}\sigma\tau)\) of the other. The trivial
character is unchanged, so both blocks have a trivial composition
factor. Write \(M_1^{\mathrm{ss}}=\mathbf1\oplus\psi_1\) and
\(M_2^{\mathrm{ss}}=\mathbf1\oplus\psi_2\). The relation above gives
\[
 \mathbf1\oplus\psi_1\simeq\overline\eta\oplus\overline\eta\psi_2.
\]
Since \(\overline\eta\neq\mathbf1\), we get
\(\overline\eta\psi_2=\mathbf1\), so \(\psi_2=\overline\eta\), and then
\(\psi_1=\overline\eta\). Hence
\[
 M_1^{\mathrm{ss}}\simeq M_2^{\mathrm{ss}}\simeq\mathbf1\oplus\overline\eta.
\]

\emph{The determinant contradiction.} On \(G_K\), the determinant of
\(\overline V\) is \(\det M_1\cdot\det M_2=\overline\eta^2=1\). By
fact (f) of Section~\ref{sec:inputs}, \(\chi_{11}^2|_{G_K}=1\). This
puts \(\Q(\zeta_{11}+\zeta_{11}^{-1})\), a field of degree \(5\), inside
the quadratic field \(K\). This is a contradiction.

We conclude that \(\Gamma\neq D_4\). No invariant maximal order was
assumed in this section.

\section{Endomorphism field of degree at most 2}\label{sec:small}

Assume case (S), let \(P\in A(\Q)\) have order \(11\), and suppose
that \(\Gamma\in\{C_1,C_2\}\).

If \(\Gamma=C_1\), then \(L=\Q\). Choose an isogeny
\(E^2\to A_{\overline{\Q}}\), and let \(B\) be the image of
\(E\times0\). By Lemma~\ref{lem:subvarieties}(a), \(B\) is an elliptic
curve over \(\Q\) contained in \(A\). This contradicts
Section~\ref{sec:subvarieties}.

If \(\Gamma=C_2\), then \(L\) is quadratic. Let \(B\) be as before. By
Lemma~\ref{lem:subvarieties}(a), it is an elliptic curve over \(L\)
contained in \(A_L\). Put \(C=A_L/B\), an elliptic curve over \(L\).
For \(\tau\in G_{\Q}\), the conjugates \(B^\tau\) and \(C^\tau\) are an
elliptic subvariety and an elliptic quotient of
\(A_{\overline{\Q}}\), because \(A\) is defined over \(\Q\). By
Section~\ref{sec:subvarieties}, all of them are isogenous to \(E\) over
\(\overline{\Q}\). Hence \(B\) and \(C\) are \(\Q\)-curves over the
quadratic field \(L\). The point \(P\) either lies in \(B(L)\) or maps
to a point of order \(11\) in \(C(L)\). Both possibilities contradict
the theorem of Le Fourn and Najman.

Together with Sections~\ref{sec:rotations} and~\ref{sec:klein}, this
proves that \(A(\Q)[11]=0\) in case (S).

\begin{remark}
For the statement about points of order \(22\), the input (c) can be
replaced here by Kenku--Momose. Let \(B'\) be the image of
\(0\times E\). It is an elliptic curve over \(L\), with \(B+B'=A_L\)
and \(B\cap B'\) finite. The quotient map therefore restricts to an
\(L\)-isogeny \(B'\to C\). By Lemma~\ref{lem:subvarieties}(b), every
isogeny \(B\to B'\) over \(\overline{\Q}\) is defined over \(L\). Hence
\(B\) and \(C\) lie in one \(L\)-isogeny class. Given a rational point
of order \(22\), its multiples of orders \(2\) and \(11\) give
\(L\)-rational points of these orders on curves in this class.
Lemma~\ref{lem:coprime} then gives an elliptic curve over the quadratic
field \(L\) with a rational point of order \(22\), contradicting input
(b). The input (c) is used only for the sharper statement that a
non-CM geometric square has no rational point of order \(11\).
\end{remark}

\section{The CM square}\label{sec:cm}

Assume case (CM), so that \(D\simeq M_2(M)\), with centre \(M\).

\emph{The degree \([L:\Q]\) is prime to \(5\).} Suppose not. By
Cauchy's theorem, \(\Gamma\) contains an element \(\gamma\) of order
\(5\). It acts on the centre \(M\) by a field automorphism. Since
\(\Aut(M/\Q)\) has order \(2\), this action is trivial. Thus
\(\gamma\) acts on \(D\) by an \(M\)-algebra automorphism. By the
Skolem--Noether theorem, it is conjugation by some
\(R\in\GL_2(M)\). Since \(\Gamma\) acts faithfully, \([R]\) has order
\(5\) in \(\PGL_2(M)\). As in the proof of
Lemma~\ref{lem:finite-subgroups}, the eigenvalues of \(R\) are distinct,
and their ratio \(\zeta\) is a primitive fifth root of unity. Hence
\[
 t=\frac{(\operatorname{tr}R)^2}{\det R}=2+\zeta+\zeta^{-1}\in M.
\]
Since \(\zeta+\zeta^{-1}\) is a root of \(X^2+X-1\), the element \(t\)
is a root of \(X^2-3X+1\), whose discriminant is \(5\). Thus
\(\sqrt5\in M\), and \(M=\Q(\sqrt5)\). This contradicts the fact that
\(M\) is imaginary quadratic.

\emph{The field \(L\) contains the relevant ring class fields.} Choose
an isogeny \(E^2\to A_{\overline{\Q}}\), and let \(B\) be the image of
\(E\times0\). By Lemma~\ref{lem:subvarieties}, \(B\) is an elliptic
curve over \(L\), and all its geometric endomorphisms are defined over
\(L\). Thus \(\End(B_L)\otimes\Q\simeq M\). The action of
\(\End(B_L)\) on the one-dimensional \(L\)-vector space \(\Lie(B)\)
gives a ring homomorphism \(\End(B_L)\otimes\Q\to L\). It is injective
because \(M\) is a field. Since \(M/\Q\) is Galois, it follows that
\(M\subseteq L\). Put \(C=A_L/B\), an elliptic curve over \(L\). It is
geometrically isogenous to \(E\), so it also has complex
multiplication by an order in \(M\). For \(Z\in\{B,C\}\), we have
\(j(Z)\in L\) and \(M\subseteq L\), so
\[
 M(j(Z))\subseteq L.
\]

\emph{Conclusion.} Suppose that \(P\in A(\Q)\) has order \(11\). Either
\(P\in B(L)\), or \(P\) maps to a point of order \(11\) in \(C(L)\). In
either case, some \(Z\in\{B,C\}\) has an \(L\)-rational point of order
\(11\). Input (d) gives \(5\mid[L:M(j(Z))]\), and hence
\(5\mid[L:\Q]\). This contradicts the first step. Hence
\(A(\Q)[11]=0\) in case (CM).

This step cannot be replaced by Clark's list of torsion orders for
abelian surfaces with everywhere potentially good reduction
\cite[Theorem~5]{Clark}, because that list contains \(22\). The divisibility in input (d) holds for
every imaginary quadratic order, including orders whose conductor is
divisible by \(11\).

\section{Completion of the proof}\label{sec:completion}

Let \(A/\Q\) be geometrically split, and suppose that \(P\in A(\Q)\)
has order \(22\). Put \(P_2=11P\) and \(P_{11}=2P\), of orders \(2\)
and \(11\). By Section~\ref{sec:subvarieties}, \(A\) is simple over
\(\Q\). By Sections~\ref{sec:rotations}--\ref{sec:cm}, \(A\) is of type
(N). These two steps prove the second assertion of
Theorem~\ref{thm:main}, since they use only \(P_{11}\).

\emph{The two elliptic subvarieties.} Choose an isogeny
\(E_1\times E_2\to A_{\overline{\Q}}\), and let \(B_i\) be the image of
\(E_i\). Then \(B_i\) is isogenous to \(E_i\), and \(B_1\cap B_2\) is
finite. Let \(B\) be any elliptic subvariety of
\(A_{\overline{\Q}}\). By Poincar\'e's theorem, \(B\) is isogenous to
\(E_1\) or to \(E_2\). The quotient \(A_{\overline{\Q}}/B_1\) is an
elliptic curve isogenous to \(B_2\), hence to \(E_2\). If \(B\) is
isogenous to \(E_1\), the composite
\(B\to A_{\overline{\Q}}\to A_{\overline{\Q}}/B_1\) is zero, because
\(E_1\) and \(E_2\) are not isogenous. Thus \(B\subseteq B_1\), and
\(B=B_1\). Similarly, \(B=B_2\) if \(B\) is isogenous to \(E_2\). Hence
\(B_1,B_2\) are the only elliptic subvarieties of
\(A_{\overline{\Q}}\).

The group \(G_{\Q}\) permutes \(\{B_1,B_2\}\). Neither curve is defined
over \(\Q\), since \(A\) is simple over \(\Q\). Hence the stabilizer of
\(B_1\) has index \(2\). It equals \(G_K\) for a quadratic field \(K\).
Let \(\sigma\) be the nontrivial automorphism of \(K\). Then
\(E=B_1\) is defined over \(K\), and \(B_2=E^\sigma\).

\emph{One isogeny class.} Put \(F=A_K/E\). The quotient map restricts
to a \(K\)-isogeny \(E^\sigma\to F\), because \(E\cap E^\sigma\) is
finite. Applying \(\sigma\), we obtain a \(K\)-isogeny
\(E\to F^\sigma\).

For each \(\ell\in\{2,11\}\), the exact sequence
\[
 0\longrightarrow E\longrightarrow A_K\xrightarrow{\ \pi\ }F
 \longrightarrow0
\]
shows the following. Either \(P_\ell\in E(K)\), or \(\pi(P_\ell)\) is a
point of order \(\ell\) in \(F(K)\). In the second case,
\(\pi^\sigma(P_\ell)=\sigma(\pi(P_\ell))\) is a point of order \(\ell\)
in \(F^\sigma(K)\), since \(P_\ell\) is rational. Thus the
\(K\)-isogeny class of \(E\) contains curves with \(K\)-rational points
of orders \(2\) and \(11\). Lemma~\ref{lem:coprime} produces an
elliptic curve over the quadratic field \(K\) with a \(K\)-rational
point of order \(22\). This contradicts input (b), and proves
Theorem~\ref{thm:main}.
\hfill\(\square\)

\section{Consequence for the census}

The groups
\href{https://genus-2-torsion.github.io/group.html?g=22}{\([22]\)} and
\href{https://genus-2-torsion.github.io/group.html?g=2.22}{\([2,22]\)}
are impossible in the class ``geometrically split, simple over
\(\Q\)''. Both groups were already impossible in the class ``split over
\(\Q\)'', by Mazur's theorem applied to an elliptic subvariety and its
elliptic quotient. Hence both geometrically split example columns are
impossible for these two groups, and their combined geometrically split
``\(\infty\)?'' entry becomes \(\emptyset\) (impossible). No other
group in the census contains an element of order \(22\).

The point of order \(2\) cannot be omitted. The census group
\href{https://genus-2-torsion.github.io/group.html?g=11}{\([11]\)} has
a geometrically split example that is simple over \(\Q\), namely
\href{https://genus-2-torsion.github.io/curve.html?id=11.b}{11.b}. By
Theorem~\ref{thm:main}, its two geometric elliptic factors are not
isogenous over \(\overline{\Q}\).

The argument applies to every geometric splitting degree, including
degrees divisible by \(2\) or \(11\). It uses no maximal integral
endomorphism ring and no principal polarization.

\begin{thebibliography}{9}
\bibitem{BC}
A. Bourdon and P. L. Clark,
\emph{Torsion points and isogenies on CM elliptic curves},
J. London Math. Soc. (2) \textbf{102} (2020), 580--622,
Theorem~4.1(a), (d), (e) (numbering of arXiv:1906.07121v1).
\href{https://arxiv.org/pdf/1906.07121}{arXiv:1906.07121}.

\bibitem{Clark}
P. L. Clark,
\emph{Bounds for torsion on abelian varieties with integral moduli},
arXiv:math/0407264v2 (2004), Theorem~5.
\href{https://arxiv.org/pdf/math/0407264}{arXiv:math/0407264}.
The same list is Theorem~1.4 of P. L. Clark and X. Xarles, \emph{Local bounds for torsion points on abelian varieties}, Canad. J. Math. \textbf{60} (2008), 532--555, \href{https://doi.org/10.4153/CJM-2008-026-x}{doi:10.4153/CJM-2008-026-x}.

\bibitem{KM}
M. A. Kenku and F. Momose,
\emph{Torsion points on elliptic curves defined over quadratic fields},
Nagoya Math. J. \textbf{109} (1988), 125--149, Theorem~(0.1).
\href{https://doi.org/10.1017/S0027763000002816}{doi:10.1017/S0027763000002816}.

\bibitem{LFN}
S. Le Fourn and F. Najman,
\emph{Torsion of \(\Q\)-curves over quadratic fields},
Math. Res. Lett. \textbf{27} (2020), 209--225, Theorem~1.1,
Proposition~3.2, and Theorem~2.1 (numbering of arXiv:1803.03460v3).
\href{https://arxiv.org/pdf/1803.03460}{arXiv:1803.03460}.

\bibitem{Mazur}
B. Mazur,
\emph{Modular curves and the Eisenstein ideal},
Publ. Math. IH\'ES \textbf{47} (1977), 33--186, Theorem~(8).
\href{https://www.numdam.org/item/PMIHES_1977__47__33_0/}{Numdam}.
\end{thebibliography}

\end{document}
